\begin{figure*}[!htb]
\small
\begin{tcolorbox}[
    colback=white, 
    colframe=blue2,
    boxrule=0.5mm,
    arc=0mm,
    outer arc=0mm,
    title=Potential Tools: Web - Generation Prompt
]
Task Description:\\
You need to go ahead and generate samples of tool usage with as much diversity as possible, based on the current dimension, the requirements requirements of the task, and the sample references provided.\\

Current dimension:\\
Application Environment Dimension - Special Environment Constraints: Web\\

Task Requirement:\\
- Generate new samples taking into account as much as possible the diversity of user query and available tools.
- Use "// Raw Input //" and "// Raw output //" to label the inputs and outputs of the samples.
- The samples need to consider as much as possible a scenario that could happen in reality, and multiple tools need to be available in the raw input. In the current dimension, the task description is in a web-related environment, including but not limited to: shopping price comparison, web content retrieval, web data collection, website usability evaluation, website security evaluation, website optimization, etc.\\

At build time, you need to:\\
Give an overall task first, then express the specific task requirements separately, and make sure that they are as relevant to the web environment as possible. In the list of available tools, make sure that you provide tools that can accomplish these requirements, while including a small number of tools that do not directly help the task. Next, repeat certain requirements once with different expressions to ensure diversity in task complexity and tools. Finally, complete the corresponding task planning process in the output based on these requirements.\\

Sample Reference:\\
==== Reference Sample Start ====\\
// Raw Input //\\
\colorbox{blue2!50}{\texttt{[sample\_input]}}\\

// Raw output //\\
\colorbox{blue2!50}{\texttt{[sample\_outputl]}}\\
==== Reference Sample End ====\\

Now, please combine the reference sample and the corresponding build steps, try to select tasks from different domains and scenarios, and use different types of tools to make the new sample different from the reference sample in terms of task content and context. We start generating new single samples corresponding to the dimensions.\\
Special Environment Constraints: Web generation sample (one):\\
==== Sample Start ====
\end{tcolorbox}
\caption{Generation prompt for potential tools scenario, Web subtask,}
\label{fig:web-prompt}
\end{figure*}